\section{Discussion and limitations}
\label{sec:limitations}
Next, we briefly discuss limitations and directions for improving the capabilities in \sys.

\para{Improve \comprehensiveness}
 We saw that \sys was unable to obtain {\em  all} critical assets in some exercises.
 In some cases, the LLM planner  obtained a single critical asset and then stopped.
In many trials, we observed the  LLM planner  could have queried the \sys attack graph service to identify that there were additional paths to explore, but did not do so.
We speculate this could be because LLMs may not have much training data for red teaming multi-host networks  using attack graphs.
An interesting direction for future work is to improve coverage (e.g., further train and fine-tune LLMs to  better leverage the attack graph service).

\para{Reducing  failure scenarios}
We saw \sys was unable to succeed  in 3 environments.
 On further analysis, we found that these settings required both external scans (e.g., to identify vulnerable web servers) and internal scans (e.g., to identify a vulnerable database management server).
 At a high level, the current LLM-assisted workflow seems to lack an   understanding  that scanning from different network locations could have different results.
We hypothesize  that this can be addressed by improving  the
\sys  attack graph service   to reason about network  segments and access control (e.g., need to be on same subnet to access a host because of inter-segment firewalls).
We believe that extending the attack graph service to reason about fine-grained access control could further improve \acquisition and also   help improve \comprehensiveness.

\para{Extending \sys to handle 0-days}
In this paper, we scoped our exploration to    consider exercises with known vulnerabilities.
Since \sys is extensible, future versions could support advanced 0-day  specific task agents to further improve red team effectiveness~\cite{wang2025cybergym}.

\para{Environment realism}
In general,  enterprise network details are  considered sensitive  and there is little public information.
\benchmark is our best effort attempt using a variety of public sources and prior reports  to design realistic environments~\cite{equifax_report, colonial_pipeline_techtarget, ciscoEnterpriseNetwork, ferguson2021_deception_psychology}.
An interesting direction of future work is to extend \benchmark and use \sys on a broader range of real  (possibly proprietary) enterprise settings at scale.

\para{Adding defenders in the loop}
As a first step toward understanding the feasibility of LLM-assisted red teaming  in multi-host network  settings, we  evaluate \sys\ in environments without defenders.
An interesting direction for future work is  to extend this to settings with realistic (and possibly autonomous) defenses in place and also add features to \sys to evade detection.


\para{Memorization}  A concern with LLMs is the memorization of training data.
Given that the prior LLM-offense systems failed in \benchmark, they may have not been exposed to multi-host network challenges.
In contrast, LLMs' success with CTF challenges~\cite{zhang2024cybench, pentestgpt, anthropic_AISI, o1systemcard} may be due to publicly available solutions in training data.
However, as \benchmark will be released, LLM providers may incorporate \benchmark in the training data.
Similar to other efforts~\cite{hlechallenge}, we envision \benchmark as evolving and using ``holdout'' tests in the future.
